\section{Tổng kết Giai đoạn 1}
    Thông qua quá trình thực hiện Đồ án Chuyên ngành (Giai đoạn 1), nhóm đã hoàn thành việc nghiên cứu các khái niệm nền tảng, làm rõ vị trí, vai trò chiến lược và cấu trúc hoạt động của IoT Gateway trong hạ tầng Internet vạn vật. Việc phân tích các công nghệ kết nối, giao thức truyền dữ liệu phổ biến cùng các thiết bị khả cấu hình hiện có trên thị trường đã tạo cơ sở vững chắc để xây dựng đặc tả chi tiết và thiết kế kiến trúc hệ thống phù hợp. Nhóm đã hiện thực hóa thành công nguyên mẫu (prototype) IoT Gateway dựa trên kiến trúc Dual-MCU (Master-Slave) sử dụng hai vi điều khiển ESP32-S3 vận hành trên nền tảng FreeRTOS. Kết quả triển khai thực tế cho thấy sản phẩm đã đạt được các mục tiêu cốt lõi về tính mô-đun hóa với các khe cắm mở rộng (WAN, LAN1, LAN2) tương thích hoàn toàn với các Module Adapter, cho phép thay đổi linh hoạt các chuẩn kết nối vật lý mà không cần thiết kế lại bo mạch chủ. Đồng thời, khả năng khả cấu hình của thiết bị được khẳng định qua ứng dụng trên PC, giúp người dùng tinh chỉnh các tham số vận hành như tài khoản MQTT hay cấu hình mạng Wi-Fi/LTE một cách dễ dàng mà không cần can thiệp vào mã nguồn. Hệ thống vận hành ổn định trên bo mạch in 6 lớp, đảm bảo toàn vẹn tín hiệu và xử lý chính xác dữ liệu từ các cảm biến sử dụng giao thức LoRa, Zigbee, CAN, RS-485. Những kết quả này cùng các cơ chế bảo vệ, tự động kết nối lại và lưu trữ đệm trên thẻ nhớ microSD đã chứng minh tính khả thi và độ tin cậy của giải pháp, tạo nền tảng vững chắc cho việc hoàn thiện sản phẩm ở giai đoạn tiếp theo.